\section*{Chapter \ref{ch:pl:security} --- Security}

\begin{solution}{exo:pl:security:1}
Because harmless roles combine into toxic people. In the chapter's policy three quantitative researchers who were also risk managers could
propose parameters and approve them, a combination neither role held alone.
\end{solution}

\begin{solution}{exo:pl:security:2}
Half the rotation period on average: 45 days under 90-day rotation, three and a half days under weekly rotation.
\end{solution}

\begin{solution}{exo:pl:security:3}
Read and trade --- which can still lose money, for example by trading against the firm's own positions --- but not withdraw: the key's scope
excludes it, and even a withdrawal key can send only to allow-listed addresses. Forged or replayed requests are refused.
\end{solution}

\begin{solution}{exo:pl:security:4}
The baseline is the median of the last thirty days; after thirty days of copying, the copying days are the baseline, the person's usual level
has risen to include them, and the daily scores fall back to zero.
\end{solution}

\begin{solution}{exo:pl:security:5}
A day of 30\% more reading is about one standard deviation of a person's daily noise; a threshold high enough to allow only one false alarm a
month per hundred people sits above four. Any single day of the copy is ordinary; only the sum of many is not.
\end{solution}

\begin{solution}{exo:pl:security:6}
A check of what was actually to be signed, independent of the screen that prepared it: the transaction decoded on the signing device or a second
system, or a policy engine that refuses a transaction type (a change to the wallet's own contract, a new destination) whatever the signatures.
\end{solution}

\begin{solution}{exo:pl:security:7}
With the thresholds recalibrated, the lagged CUSUM flags all 20 insiders at +30\% a day, in a median of 20.5 days instead of 22 (19 of 20). A
smaller reference value lets smaller daily excesses accumulate: it is tuned to a smaller shift, at the price of a higher threshold for the same
false-alarm rate.
\end{solution}

\begin{solution}{exo:pl:security:8}
The firewall is one control, not a reason for trust: a compromised laptop, a stolen credential or a malicious insider is already inside. Every
request should be authenticated and authorised on its own, by identity and device, with paths between zones stated and minimal.
\end{solution}

\begin{solution}{pb:pl:security:1}
\begin{enumerate}
\item Money that moves on a signature, from criminals and states; information whose value is its secrecy, from competitors and departing staff.
\item No submitting orders and amending trades; no proposing and approving parameters; no merging code and deploying it; no creating venue
keys and withdrawing; no creating and releasing payments.
\item Developer, head of desk, operations, site reliability and trader; 60 of the 100 people.
\item Split the five roles, named a release manager, removed the three second roles: no toxic role and no toxic person.
\item In the policy's own \omterm{def:pl:testing-and-continuous-delivery:ci}{continuous integration}, failing any change that creates a combination.
\item Leases, access logs, rotation and revocation, and no copy in code or configuration.
\item A leaked secret is valid only until the next rotation: half the period on average.
\item A known key, a valid signature, a fresh timestamp, an action within the key's scope, and an allow-listed withdrawal address.
\item The trading key: reads and trades; the withdrawal key: withdrawals to the custodian only.
\item It keeps private keys from ever leaving it; it does not check what it is asked to sign.
\item Robust scores against each person's own 30-day baseline, trailing or lagged by 60 days, used alone or accumulated in a CUSUM; each threshold
set on clean people for one false alarm a month per hundred.
\item Single day, trailing: 1 of 20; CUSUM, trailing: 5; single day, lagged: 8; CUSUM, lagged: 19, in a median of 22 days.
\item It adds small daily excesses that no single day shows.
\item A trailing baseline absorbs the copying within a month; a lagged one keeps it out of its own reference.
\item On day 131, eleven days after the copying began.
\item \textbf{Named result.} At one false alarm a month per hundred people, an insider reading 30\% more strategy code than usual every day is
flagged in a median of 22 days (19 of 20) by a CUSUM on a lagged baseline, and 50\% more in 12 days (20 of 20); a single-day threshold on a
trailing baseline flags one of twenty. The policy held 5 toxic roles and 60 toxic people before the fix, none after.
\item Peer comparison within the role and features of novelty (files never read before).
\item Trading, research, corporate and outside zones, with orders only to venues, fills to research, parameters back only through review, and
no path from laptops or the internet to trading.
\item Verify what is signed independently of what prepares it, and let policy refuse what no signer should approve.
\item Who and what may do what, and whether what they see is what they do.
\end{enumerate}
\end{solution}

\begin{solution}{iq:pl:security:1}
In a secrets vault, issued to the strategy's service identity under a lease at start-up, rotated on a schedule, never in code, images or
configuration, and every access logged.
\iqlookfor{Vault, leases, rotation.}
\end{solution}

\begin{solution}{iq:pl:security:2}
Read and trade only, never withdraw; restricted to the firm's IP addresses; a separate key per strategy or server; withdrawals, if needed at
all, on a separate key limited to allow-listed addresses.
\iqlookfor{Scopes and allow-lists.}
\end{solution}

\begin{solution}{iq:pl:security:3}
Log reads of protected code and data per person; score each day against the person's own lagged baseline and their peers; accumulate the
scores with a CUSUM; calibrate the threshold on clean data to a false-alarm rate the team can investigate; add data-loss controls on copies
and uploads.
\iqlookfor{Accumulation and a baseline that does not learn the theft.}
\end{solution}

\begin{solution}{iq:pl:security:4}
No implicit trust from being inside the network: every request authenticated and authorised by identity and device. Inside a trading network
it means service identities, mutual authentication between services, and paths between zones stated and minimal.
\iqlookfor{Identity per request.}
\end{solution}

\begin{solution}{iq:pl:security:5}
Roles as reviewed data with segregation-of-duties checks in CI; multi-factor, phishing-resistant authentication; a vault with leases and
rotation; scoped venue keys and allow-listed withdrawals on separate keys; treasury signing with independent transaction verification and
policy limits; network zones; access logging with calibrated insider detection; a \omterm{def:pl:security:threat}{threat model} reviewed yearly.
\iqlookfor{\omterm{def:pl:security:rbac}{Least privilege}, bounded secrets, verified signing, detection.}
\end{solution}
